% Zdroj: PDF priklady-leto-2024.pdf, fyzická str. 23-24. Značka: specializace UI-SU, UI-ZPJ. Zařazeno: S3.
\section{PCA analýza}
\szzotazka{léto 2024}{29}{PCA analýza}

\noindent\textbf{Zadání.}\par
Mějme data se dvěma atributy a ze třech tříd zobrazená na obrázku níže.

\begin{center}
\includegraphics[width=0.7\linewidth]{pca-leto2024.png}
\end{center}

\begin{enumerate}
  \item Vysvětlete, jak funguje PCA analýza. Jakým způsobem můžeme spočítat hlavní komponenty?
  \item Bez počítání odhadněte, jaké budou směry dvou hlavních komponent pro data na obrázku výše. Zakreslete je do
    obrázku.
  \item Uvažujme klasifikátor založený na metodě $k$ nejbližších sousedů pro vhodné $k$ a předpokládejme, že jsme v rámci
    předzpracování dat použili PCA analýzu pro snížení počtu dimenzí na 1. Jaká bude přesnost tohoto klasifikátoru v
    porovnání s přesností stejného klasifikátoru trénovaného na původních datech? Vysvětlete.
\end{enumerate}

\medskip
\noindent\textbf{Řešení.} \textit{(V~PDF bez náčrtu řešení; vlastní vzorové řešení.)}\par
\begin{enumerate}[nosep]
  \item PCA: data se vycentrují (odečte se průměr), spočte se \emph{kovarianční matice} a~její \emph{vlastní vektory} -- to jsou hlavní komponenty, seřazené sestupně podle vlastních čísel (každé udává rozptyl zachycený v~daném směru). Redukce dimenze = projekce na prvních pár komponent.
  \item První hlavní komponenta míří ve směru \emph{největšího rozptylu} dat (podél protáhlého mraku, zhruba po antidiagonále), druhá kolmo na ni.
  \item Přesnost se spíš \emph{zhorší}. PCA je metoda \emph{bez učitele} -- maximalizuje rozptyl, nikoli oddělitelnost tříd. Na těchto datech leží rozdíl mezi třídami převážně ve směru \emph{kolmém} na 1.~komponentu, takže projekce na $1$D třídy překryje a~$k$-NN ztratí informaci potřebnou k~jejich rozlišení.
\end{enumerate}
